% TOP_PROBLEM: {"id": 1181, "title": "The Pierce-Birkhoff Conjecture", "category_id": 4, "category": {"id": 4, "name": "algebra", "display_name": "Algebra", "description": "Group theory, ring theory, field theory, and algebraic structures.", "slug": "algebra", "order_index": 4, "created_at": "2026-07-31T15:26:25.671Z"}, "status": "open"}
% FIRST_POSTED: null
% ENTRY_KIND: statement_only
% Imported from: batch04.tex; UnsolvedMath ID 1181
\documentclass[11pt]{article}
\usepackage[margin=25mm]{geometry}
\usepackage{fontspec}
\setmainfont{Latin Modern Roman}
\usepackage{amsmath,amssymb,mathtools}
\usepackage{unicode-math}
\setmathfont{Latin Modern Math}
\usepackage{xurl,hyperref}
\hypersetup{colorlinks=true,urlcolor=blue}
\setcounter{secnumdepth}{0}
\setlength{\parindent}{0pt}
\setlength{\parskip}{5pt}
\setlength{\emergencystretch}{3em}
\newcommand{\sref}[2]{\hyperlink{source:#1:#2}{[S#2]}}
\newcommand{\cataloguescope}[1]{\paragraph{Recorded target.}#1}
\begin{document}
% UnsolvedMath ID 1181; source catalogue code ALG-021
\setcounter{section}{1180}
\section{The Pierce-Birkhoff Conjecture}\label{1181}
{\small\sffamily UnsolvedMath ID 1181 · Algebra}
\subsection{Definitions and mathematical statement}

\paragraph{Shared notation.}$\mathbb N=\{1,2,\ldots\}$, and $\mathbb Z,\mathbb Q,\mathbb R,\mathbb C$ denote the integers, rationals, reals and complex numbers. Any other conventions are specified locally.


A semialgebraic subset of \(\mathbb R^n\) is obtained from finitely
many polynomial equalities and inequalities by finitely many Boolean
operations.  A function \(f:\mathbb R^n\to\mathbb R\) is piecewise
polynomial in this problem if there are closed semialgebraic sets
\(C_1,\ldots,C_r\) covering \(\mathbb R^n\) and polynomials
\(p_1,\ldots,p_r\in\mathbb R[X_1,\ldots,X_n]\) such that
\(f|_{C_i}=p_i|_{C_i}\).  In particular, \(f\) is continuous;
arbitrary discontinuous piecewise formulas are excluded.
\paragraph{Statement recorded in the supplied source.}
For every \(n\ge1\) and every such continuous piecewise-polynomial
function \(f:\mathbb R^n\to\mathbb R\), do there exist a positive
integer \(s\), positive integers \(t_1,\ldots,t_s\), and polynomials
\(q_{ij}\in\mathbb R[X_1,\ldots,X_n]\) such that, for every
\(x\in\mathbb R^n\),
\[
 f(x)=\max_{1\le i\le s}\ \min_{1\le j\le t_i}q_{ij}(x)?
\] \sref{1181}{2}
\subsection{Short English statement}
Can every continuous piecewise-polynomial real function be expressed as a finite maximum of finite minima of polynomials?
\subsection{Sources}
\begin{description}
\item[\hypertarget{source:1181:1}{S1}] ULAM.AI and the UnsolvedMath Contributors, \emph{UnsolvedMath: A Curated Collection of Open Mathematics Problems}, record 1181 (ALG-021). CC BY 4.0. \url{https://huggingface.co/datasets/ulamai/UnsolvedMath}.
\item[\hypertarget{source:1181:2}{S2}] Fran\c{c}ois Lucas, James Madden, Daniel Schaub and Mark Spivakovsky, \emph{Approximate roots of a valuation and the Pierce--Birkhoff Conjecture}, arXiv:1003.1188, Section 0.1, Definition 0.1.1 and Conjecture 0.1.2. \url{https://arxiv.org/abs/1003.1188}.
\end{description}
\label{end:1181}
\end{document}
